\section{Корректность исчисления предикатов}%
\index{Корректность исчисления предикатов|defin}%
\index{Исчисление!предикатов!корректность|defin}%
        \label{predicate-correctness}

\begin{theorem}
        %
Всякая выводимая в исчислении предикатов формула является общезначимой.
        %
\end{theorem}

\begin{proof}
        %
Для исчисления высказываний проверка корректности была
тривиальной\т надо было по таблице проверить, что все аксиомы
(1)\ч (11) являются тавтологиями. С этими аксиомами и сейчас нет
проблем. Но в двух следующих аксиомах есть ограничение на
корректность подстановки,%
\index{Корректная подстановка}
без которого они могут не быть
общезначимыми. Естественно, это ограничение должно быть
использовано и в доказательстве корректности, и это потребует
довольно скучных рассуждений\т тем более скучных, что сам факт
кажется ясным и так. Тем не менее такие рассуждения надо уметь
проводить, так что мы ничего пропускать не будем.

Итак, пусть фиксирована сигнатура $\sigma$, а также некоторая
интерпретация этой сигнатуры. Всюду далее, говоря о термах и
формулах, мы имеем в виду термы и формулы этой сигнатуры, а
говоря об их значениях, имеем в виду значения в этой
интерпретации.

\textsf{Лемма 1}. Пусть $u$ и $t$\т термы, а $\xi$\т
переменная. Тогда
        $$
   [u (t/\xi)](\pi) = [u](\pi + (\xi\mapsto [t](\pi)))
        $$
для произвольной оценки $\pi$.

Напомним обозначения: в левой части мы подставляем $t$ вместо $\xi$
в терм $u$, и берём значение получившегося терма на оценке $\pi$.
В правой части стоит значение терма~$u$ на оценке, которая получится
из $\pi$, если значение переменной $\xi$ изменить и считать равным
значению терма $t$ на оценке $\pi$.

В сущности, это утверждение совершенно тривиально: оно говорит,
например, что значение $\sin(\cos(x))$ при $x=2$ равно значению
$\sin(y)$ при $y=\cos(2)$. Но раз уж мы взялись всё доказывать
формально, докажем его индукцией по построению терма $u$. Если
терм $u$ есть переменная, отличная от $\xi$, то ни подстановка,
ни изменение оценки не сказываются на значении терма $u$. Для
случая~$u=\xi$ получаем $[t](\pi)$ слева и справа. Если терм
получается из других термов применением функционального символа,
то подстановка выполняется отдельно в каждом из этих термов, так
что искомое равенство также сохраняется. Лемма~1 доказана.

Аналогичное утверждение для формул таково:

\textsf{Лемма 2}. Пусть $\varphi$\т формула, $t$ \т терм, а
        \label{trouble-lemma}%
$\xi$\т переменная, причём подстановка $t$ вместо $\xi$ в
формулу $\varphi$ корректна. Тогда
        $$
   [\varphi(t/\xi)](\pi) = [\varphi](\pi + (\xi\mapsto [t](\pi)))
        $$
для произвольной оценки $\pi$.

Поясним смысл этой леммы на примере. Пусть $\xi$ является
единственным параметром формулы~$\varphi$, а $c$\т константа.
Тогда формула $\varphi(c/\xi)$ замкнута; лемма утверждает, что
её истинность равносильна истинности $\varphi$ на оценке, при
которой значение переменной $\xi$ есть элемент интерпретации,
соответствующий константе~$c$.

Доказательство леммы проведём индукцией по построению
формулы~$\varphi$. Для атомарных формул это утверждение является
прямым следствием леммы~1. Кроме того, из определения
истинностного значения формулы и из определения подстановки
ясно, что если утверждение леммы~2 верно для двух формул
$\varphi_1$ и $\varphi_2$, то оно верно для их любой их
логической комбинации (конъюнкции, дизъюнкции и импликации);
аналогично для отрицания.

Единственный нетривиальный случай\т формула, начинающаяся с
квантора. Здесь наши определения вступают в игру.
Пусть $\varphi$ имеет вид $\forall \eta\, \psi$. Есть два
принципиально разных случая: либо $\xi$ является параметром
формулы $\varphi$ (\те формулы
$\forall \eta\,\psi$), либо нет. Во втором случае
$\varphi(t/\xi)$ совпадает с $\varphi$, а изменение значения
переменной~$\xi$ в оценке~$\pi$ не влияет на значение формулы
$\varphi$, так что всё сходится. Осталось разобрать случай, когда
$\xi$ является параметром формулы~$\forall \eta\,\psi$ (отсюда
следует, что $\xi$ не совпадает с $\eta$). По определению
корректной подстановки, в этом случае переменная $\eta$ не
входит в терм $t$ и подстановка $\psi(t/\xi)$ корректна. Тогда
        %
\begin{multline*}
  [(\forall \eta\, \psi)(t/\xi)](\pi)=
       [\forall \eta\, (\psi(t/\xi))](\pi) =
   \wedge_m [\psi(t/\xi)](\pi+(\eta\mapsto m)) = \\
   =\wedge_m [\psi](\pi+(\eta\mapsto m)+
     (\xi\mapsto[t](\pi+(\eta\mapsto m)))).
\end{multline*}
        %
Мы воспользовались определением подстановки, определением истинности
($\wedge_m$ означает конъюнкцию по всем элементам из носителя
интерпретации) и предположением индукции для формулы $\psi$.
Теперь надо заметить, что переменная $\eta$ не входит в $t$ по
предположению корректности, и потому значение терма $t$ не изменится,
если заменить $\pi+(\eta\mapsto m)$ на $\pi$. Далее, $\xi$ и $\eta$
различны, поэтому два изменения в $\pi$ можно переставить местами.
Используя эти соображения, можно продолжить цепочку равенств:
        %
\begin{align*}
        &=\wedge_m [\psi](\pi+(\xi\mapsto [t](\pi))+(\eta\mapsto m)) = \\
        &=[\forall \eta\,\psi](\pi+(\xi\mapsto [t](\pi))) = \\
        &=[\varphi](\pi+(\xi\mapsto [t](\pi))),
\end{align*}
        %
что и требовалось. Случай формулы вида $\exists\xi\,\psi$
разбирается аналогично, надо только $\wedge_m$ заменить на $\vee_m$.
Лемма~2 доказана.

Теперь уже ясно, почему формула
        $$
\forall \xi \, \varphi \to \varphi(t/\xi)
        $$
будет истинна на любой оценке $\pi$ (если подстановка
корректна). В самом деле, если левая часть импликации истинна
на~$\pi$, то $\varphi$ будет истинна на любой оценке $\pi'$,
которая отличается от $\pi$ лишь значением переменной $\xi$. В
частности, $\varphi$ будет истинна и на оценке $\pi+(\xi\mapsto
[t](\pi))$, что по только что доказанной лемме~2 означает, что
правая часть импликации истинна на $\pi$.

Общезначимость формулы
        $$
\varphi(t/\xi) \to \exists \xi\,\varphi
        $$
доказывается аналогично.

Осталось проверить, что правила вывода сохраняют общезначимость.
Для правила MP это очевидно (как и в случае исчисления высказываний).
Проверим это для правил Бернайса. Это совсем просто, так как здесь
нет речи ни о каких корректных подстановках.

Пусть, например, формула $\psi \to \varphi$ общезначима и
переменная $\xi$ не является параметром формулы $\psi$.
Проверим, что формула $\psi\to\forall\xi\,\varphi$ общезначима,
то есть истинна на любой оценке $\pi$ (в любой интерпретации). В
самом деле, пусть $\psi$ истинна на оценке $\pi$. Тогда она
истинна и на любой оценке $\pi'$, отличающейся от $\pi$ только
значением переменной $\xi$ (значение переменной~$\xi$ не влияет
на истинность $\psi$, так как~$\xi$ не является
параметром). Значит, и формула $\varphi$ истинна на любой такой
оценке $\pi'$. А~это в точности означает, что $\forall
\xi\,\varphi$ истинна на оценке $\pi$, что и требовалось.

Для второго правила Бернайса рассуждение симметрично.
Пусть формула $\varphi\to\psi$ общезначима и
переменная $\xi$ не является параметром формулы $\psi$. Покажем,
что формула $\exists \xi\,\varphi\to\psi$ общезначима. В самом
деле, пусть её левая часть истинна на некоторой оценке $\pi$. По
определению истинности формулы, начинающейся с квантора
существования, это означает, что найдётся оценка~$\pi'$, которая
отличается от $\pi$ только на переменной $\xi$, для которой
$[\varphi](\pi')$ истинно. Тогда и $[\psi](\pi')$ истинно. Но
переменная $\xi$ не является параметром формулы $\psi$, так что
$[\psi](\pi')=[\psi](\pi)$. Следовательно, формула $\psi$
истинна на оценке $\pi$, что и требовалось доказать.
        %
\end{proof}

\section{Выводы в исчислении предикатов}

\subsection*{Примеры выводимых формул}
Прежде чем доказывать теорему Гёделя о полноте
исчисления предикатов, мы должны приобрести некоторый
опыт построения выводов в этом исчислении.

\begin{itemize}
        \item
Прежде всего отметим, что возможность сослаться на теорему о
полноте исчисления высказываний и считать выводимым любой
частный случай пропозициональной тавтологии сильно облегчает
жизнь. Например, пусть мы вывели две формулы $\varphi$ и $\psi$
и хотим теперь вывести формулу $(\varphi\land\psi)$. Это просто:
заметим, что формула $(\varphi\to(\psi\to(\varphi\land\psi)))$
является частным случаем пропозициональной тавтологии (а на
самом деле и аксиомой) и дважды применяем правило MP.
        \item
Другой пример такого же рода: если формула $\varphi\hm\to\psi$ выводима,
то выводима и формула $\lnot\psi\to\lnot\varphi$, поскольку импликация
$$(\varphi\to\psi)\to(\lnot\psi\to\lnot\varphi)$$ является частным
случаем пропозициональной тавтологии.
        \item
Ещё один пример: если выводимы формулы $\varphi\to\psi$ и $\psi\to\tau$,
то выводима и формула $\varphi\to\tau$, поскольку формула
        $$
(\varphi\to\psi) \to ((\psi\to\tau)\to(\varphi\to\tau))
        $$
является частным случаем пропозициональной тавтологии.
        \item
Для произвольной формулы $\varphi$ выведем формулу
        $$
\forall x \, \varphi \to \exists x \, \varphi.
        $$
В самом деле, подстановка переменной вместо себя
всегда допустима, поэтому формулы $\forall x\,\varphi \hm\to
\varphi$ и $\varphi\hm\to\exists x\,\varphi$ являются аксиомами. Остаётся сослаться на предыдущее замечание.
        \item
Для произвольной формулы $\varphi$ выведем формулу
        $$
\exists y\, \forall x\, \varphi \to
\forall x\, \exists y\, \varphi.
        $$
Формулы $\forall x\,\varphi \to \varphi$ и $\varphi\to\exists
y\,\varphi$ являются аксиомами. С их помощью выводим
формулу $\forall x\, \varphi\hm\to \exists
y\,\varphi$. Теперь заметим, что левая часть импликации не имеет
параметра $x$, а правая часть не имеет параметра~$y$, так что
можно применить два правила Бернайса (в~любом порядке) и
добавить справа квантор~$\forall x$, а слева\т квантор $\exists
y$.
        \item
Предположим, что формула $\varphi\to\psi$ выводима, а $\xi$\т
произвольная переменная. Покажем, что в этом случае выводима
формула $\forall \xi\,\varphi\hm\to \forall\xi\,\psi$. В самом
деле, формула $\forall \xi\,\varphi\hm\to\varphi$ является
аксиомой. Далее выводим (с помощью пропозициональных тавтологий
и правила MP) формулу $\forall \xi\,\varphi \to\psi$; остаётся
воспользоваться правилом Бернайса (здесь важно, что $\xi$
не является параметром левой части).
        \item
Аналогичным образом из выводимости формулы $\varphi\to\psi$
следует выводимость формулы $\exists\xi\,\varphi\hm\to\exists\xi\,\psi$,
только надо начать с аксиомы $\psi\hm\to\exists\xi\,\psi$, затем
получить $\varphi\hm\to\exists\xi\,\psi$, а потом применить
правило Бернайса.
        \item
\label{adding-quantifiers}%
Таким образом, если формулы $\varphi$ и $\psi$ доказуемо
эквивалентны (это значит, что импликации $\varphi\hm\to\psi$ и
$\psi\hm\to\varphi$ выводимы), то формулы $\forall\xi\,\varphi$
и $\forall\xi\,\psi$ также доказуемо эквивалентны. (Аналогичное
утверждение верно и для формул $\exists\xi\,\varphi$ и
$\exists\xi\,\psi$.)

Теперь несложно доказать и более общий факт: замена подформулы
на доказуемо эквивалентную даёт доказуемо эквивалентную формулу
(см.~с.~\pageref{changing-parts}).

        \item
        \label{changing-variables-example}%
Выведем формулу $\forall x\, A(x) \to \forall y\, A(y)$ (здесь
$A$\т одноместный предикатный символ). Это несложно: начнём с
аксиомы $\forall x\, A(x) \to A(y)$, в ней левая часть не имеет
параметра $y$ и потому по правилу Бернайса из неё получается
искомая формула. Этот пример показывает, что связанные
переменные можно переименовывать, не меняя смысла формулы
(подробнее об этом мы скажем в
разделе~\ref{renaming-variables}.)

        \item
        \label{quantifier-duality}%
Выведем формулы, связывающие кванторы всеобщности и
существования:
        %
\begin{align*}
        %
\forall \xi\,\varphi &\leftrightarrow \lnot\exists\xi\,\lnot\varphi;\\
\exists \xi\,\varphi &\leftrightarrow \lnot\forall\xi\,\lnot\varphi.
        %
\end{align*}
Напомним, что $\alpha\leftrightarrow\beta$ мы считаем
сокращением для ${(\alpha\to\beta)}\hm\land{(\beta\to\alpha)}$,
так что нам надо вывести четыре формулы.

Начнём с формулы $\exists\xi\,\varphi\hm\to
\lnot\forall\xi\,\lnot\varphi$. Имея в виду правило Бернайса,
достаточно вывести формулу
$\varphi\hm\to\lnot\forall\xi\,\lnot\varphi$. Тавтология
$(B\to \lnot A)\hm\to(A\to \lnot B)$ позволяет вместо этого выводить
формулу $\forall\xi\,\lnot\varphi\hm\to\lnot\varphi$, которая
является аксиомой.

В только что выведенной формуле
$\exists\xi\,\varphi\to\lnot\forall\xi\,\lnot\varphi$ можно в
качестве $\varphi$ взять любую формулу, в том числе начинающуюся
с отрицания. Подставив $\lnot\varphi$ вместо $\varphi$, получим
        $$
\exists\xi\,\lnot\varphi\to\lnot\forall\xi\,\lnot\lnot\varphi,
        $$
где $\lnot\lnot\varphi$ доказуемо эквивалентна $\varphi$ и
потому может быть заменена на $\varphi$. После этого правило
контрапозиции (если из~$A$ следует~$\lnot B$, то из~$B$
следует~$\lnot A$) даёт
        $$
\forall \xi\,\varphi \to \lnot\exists\xi\,\lnot \varphi.
        $$

Выведем третью формулу:
$\lnot\exists\xi\,\lnot\varphi\to\forall\xi\,\varphi$. По
правилу Бернайса достаточно вывести
$\lnot\exists\xi\,\lnot\varphi\to\varphi$, что после
контрапозиции превращается в аксиому
$\lnot\varphi\to\exists\xi\,\lnot\varphi$.

Четвёртая формула получится, если заменить в третьей $\varphi$
на $\lnot\varphi$ и применить контрапозицию.
        %
\end{itemize}

\subsection*{Выводимость из посылок}

В исчислении высказываний важную роль играло понятие выводимости
из посылок и связанная с ним лемма о дедукции
(с.~\pageref{deduction-lemma}). Для исчисления предикатов
ситуация немного меняется. Если разрешить использовать посылки
наравне с аксиомами безо всяких ограничений, то утверждение,
аналогичное лемме о дедукции, будет неверным. Например, из
формулы $A(x)$ можно вывести формулу $\forall x\, A(x)$ (как мы
видели на с.~\pageref{generalization} при обсуждении правила
обобщения). Но импликация $A(x)\to\forall x\,A(x)$ не является
выводимой (поскольку не общезначима).

Поэтому мы ограничимся случаем, когда все посылки являются
замкнутыми формулами. Пусть $\Gamma$\т произвольное множество
замкнутых формул рассматриваемой нами сигнатуры~$\sigma$. (Такие
множества называют \emph{теориями}\index{Теория|defin} в
сигнатуре $\sigma$.) Говорят, что формула $A$ \emph{выводима
из~$\Gamma$},%
   \index{Выводимость!формулы из $\Gamma$|defin}%
   \index{Формула!выводимая из $\Gamma$|defin}
если её можно вывести, используя наравне с аксиомами формулы
из~$\Gamma$. Как и для исчисления высказываний, мы пишем
$\Gamma\vdash A$. Выводимые из $\Gamma$ формулы называют также
\emph{теоремами}\index{Теорема!теории|defin} теории~$\Gamma$.

\textsf{Лемма о дедукции для исчисления
предикатов.}\label{deduction-predicate}\index{Лемма!о дедукции}
Пусть $\Gamma$\т множество замкнутых формул, а $A$\т замкнутая
формула. Тогда $\Gamma \vdash (A\to B)$ тогда и только тогда,
когда $\Gamma\cup\{A\} \vdash B$.

Доказательство проходит по той же схеме, что и для исчисления
высказываний (с.~\pageref{deduction-lemma}): к формулам
$C_1,\dots,C_n$, образующим вывод $C_n=B$ из $\Gamma\cup\{A\}$,
мы приписываем посылку $A$ и дополняем полученную последовательность
        $$
(A\to C_1),\dots,(A\to C_n)
        $$
до вывода из $\Gamma$. Отличие от пропозиционального случая в том,
что в выводе могут встречаться правила Бернайса. Например,
от выводимости формулы
        $$A\to(\psi\to\varphi)$$
надо перейти к выводимости формулы
        $$A\to(\psi\to\forall\xi\,\varphi)$$
(в которой переменная $\xi$ не является параметром формулы
$\psi$). Это несложно сделать, если заметить, что в силу
пропозициональных тавтологий можно перейти от
$A\to(\psi\to\varphi)$ к $(A\land\psi)\to\varphi$, затем
применить правило Бернайса (это законно, так как переменная
$\xi$ не является параметром формулы $\psi$, а формула $A$
замкнута по предположению). Получится выводимая из~$\Gamma$
формула
        $$(A\land\psi)\to\forall\xi\,\varphi,$$
и остаётся вернуть $A$ из конъюнкции в
посылку.

Сходным образом рассматривается и второе правило Бернайса. Если
выводима формула
        $A\to(\varphi\to\psi),$
то в силу пропозициональных тавтологий выводима формула
        $\varphi\to(A\to\psi),$
к которой можно применить правило Бернайса и получить
        $\exists\xi\,\varphi\to(A\to\psi),$
после чего вернуть $A$ назад с помощью пропозициональной
тавтологии. Лемма о дедукции доказана.

Отметим теперь несколько полезных свойств выводимости из посылок.

\begin{itemize}
        \item
Если $\Gamma\vdash A$ и $\Gamma' \supset \Gamma$, то $\Gamma' \vdash A$.
(Очевидно следует из определения.)
        \item
Если $\Gamma\vdash A$, то существует конечное множество $\Gamma'
\subset \Gamma$, для которого $\Gamma' \vdash A$. (Вывод конечен
и потому может использовать лишь конечное число формул.)
        \item
Если $\Gamma$ конечно и равно
$\{\gamma_1,\dots,\gamma_n\}$, то $\Gamma\vdash A$ равносильно
выводимости (без посылок)
формулы
     $$(\gamma_1\land\ldots\land\gamma_n)\to A.$$

В самом деле, если $\{\gamma_1,\dots,\gamma_n\}\vdash A$, то
многократное применение леммы о дедукции даёт
        $$
\vdash \gamma_1\to(\gamma_2\to(\ldots(\gamma_n\to A)\ldots)),
        $$
и остаётся воспользоваться надлежащей пропозиональной
тавтологией. (В обратную сторону рассуждение также проходит без
труда.)

        \item
Комбинируя три предыдущих замечания, приходим к такому
эквивалентному определению выводимости из посылок: $\Gamma\vdash A$,
если найдутся формулы $\gamma_1,\dots,\gamma_n\in \Gamma$, для которых
        $$
\vdash \gamma_1\to(\gamma_2\to(\ldots(\gamma_n\to A)\ldots)).
        $$
Это определение имеет смысл и для формул с параметрами, так что если уж
определять выводимость из посылок с параметрами (чего обычно избегают),
        \index{Выводимость!из посылок с параметрами|defin}%
то именно так.
        %
\end{itemize}

Понятие выводимости из посылок позволяет переформулировать
теорему о корректности%
\index{Корректность исчисления предикатов}%
\index{Теорема!о корректности исчисления предикатов}
исчисления предикатов.

Говорят, что интерпретация~$M$ сигнатуры~$\sigma$ является
\emph{моделью}\index{Модель!теории|defin} теории~$\Gamma$, если
все формулы из $\Gamma$ истинны в~$M$.

\begin{theorem}[о корректности; переформулировка]
\index{Теорема!о корректности исчисления предикатов, переформулировка|defin}%
  \label{correctness-reformulation}%
        %
Все теоремы теории~$\Gamma$ истинны в любой модели~$M$
теории~$\Gamma$.
        %
\end{theorem}

\begin{proof}
        %
Если формула $A$ является теоремой теории~$\Gamma$ (\те
$\Gamma\vdash A$), найдутся формулы
$\gamma_1,\dots,\gamma_n\in\Gamma$, для которых
        $$
\vdash \gamma_1\to(\gamma_2\to(\ldots(\gamma_n\to A)\ldots)).
        $$
По теореме о корректности (в уже известной нам форме) эта
формула будет истинна во всех интерпретациях, в частности в $M$.
Поскольку $\gamma_1,\dots,\gamma_n$ истинны в~$M$, то и
формула~$A$ истинна в~$M$ (на любой оценке).
        %
\end{proof}

В следующих задачах\т и только в них\т знак $\vdash$ понимается
в описанном выше смысле (в посылках допускаются параметры).

\begin{problem}
        %
Пусть $\Gamma$\т множество произвольных (не обязательно замкнутых)
формул.
        \index{Лемма!о дедукции}%
        \textbf{(а)}~%
Пусть существует \лк вывод\пк\ некоторой
формулы $\varphi$, в котором наравне с аксиомами используются
формулы из~$\Gamma$, при этом все применения правил Бернайса
предшествуют появлению формул из~$\Gamma$. Покажите, что
$\Gamma\vdash\varphi$. Покажите, что верно и обратное
утверждение.
        \textbf{(б)}~%
Покажите, если в \лк выводе\пк\ формулы $\varphi$ наравне с
аксиомами используются формулы из $\Gamma$, но правила Бернайса
не применяются по переменным, свободным в~$\Gamma$, то
${\Gamma\vdash\varphi}$.
        %
\end{problem}

\begin{problem}
        %
Покажите, что правила Бернайса можно переписать так:
       \index{Бернайса правила}%
       \index{Правило!Бернайса}%
        $$
\frac{\Gamma, A\vdash B\mathstrut}
        {\Gamma, A \vdash \forall\xi\,B\mathstrut}
\qquad
\frac{\Gamma, B\vdash A\mathstrut}
        {\Gamma, \exists\xi\,B\vdash A\mathstrut}\ ,
        $$
где переменная $\xi$ не является параметром
формулы~$A$, а также параметром формул из~$\Gamma$.
(В первом правиле мы для симметрии выделили формулу~$A$,
хотя она ничем не отличается от формул из~$\Gamma$.)
        %
\end{problem}

\subsection*{Переменные и константы}
        %
Отметим ещё несколько простых свойств выводимости, которые
нам потребуются:

\textsf{Лемма о свежих константах}.%
\index{Лемма!о свежих константах|defin}%
\index{Константа!свежая}
Пусть выводима формула~$\varphi(c/\xi)$, где $\varphi$\т
произвольная формула, $\xi$\т переменная,~$c$\т константа, не
входящая в формулу~$\varphi$. Тогда выводима и формула
$\varphi$.
        \label{lemma-fresh-constants}%

Интуитивный смысл леммы: если мы доказали что\д то про \лк
свежую\пк\ константу $c$ (не запятнавшую себя участием в формуле
$\varphi$), то фактически мы доказали формулу $\varphi$ для
произвольных значений переменной.

Доказательство леммы. По условию существует вывод формулы
$\varphi(c/\xi)$. Возьмём \лк свежую\пк\ переменную $\eta$, не
встречающуюся в этом выводе, и всюду заменим в нём константу~$c$
на эту переменную. При этом вывод останется выводом, так как
правила обращения с переменными и константами ничем не
отличаются (кванторов по новой переменной в нём нет, так что
корректные подстановки останутся корректными и применения
правил Бернайса останутся допустимыми). Таким образом, выводима
формула~$\varphi(\eta/\xi)$.

По правилу обобщения выводима формула
$\forall\eta\,\varphi(\eta/\xi)$. Осталось применить аксиому
$\forall\eta\,\varphi(\eta/\xi)\hm\to
\varphi(\eta/\xi)(\xi/\eta)$; подстановка в правой части
корректна и даёт формулу~$\varphi$, так как сначала мы заменили
свободные вхождения $\xi$ на~$\eta$, а затем обратно (так что в
зону действия кванторов по $\xi$ они попасть не могли). Лемма
доказана.

\begin{problem}
        %
Сформулируйте и докажите аналогичную лемму для нескольких констант.
        %
\end{problem}

Аналогичное рассуждение позволяет доказать и другое утверждение,
которое нам потребуется:

\textsf{Лемма о добавлении констант}.%
\index{Лемма!о добавлении констант|defin}
Пусть формула $\varphi$
некоторой сигнатуры~$\sigma$ выводима в исчислении предикатов
расширенной сигнатуры~$\sigma'$, полученной из $\sigma$
добавлением новых констант. Тогда $\varphi$ выводима и в
исчислении предикатов сигнатуры~$\sigma$.
        \label{lemma-adding-constants}%

Доказательство. Пусть формула $\varphi$, не содержащая новых
констант, имеет вывод, в котором новые константы встречаются.
Как их оттуда удалить? Легко понять, что их можно заменить на
свежие переменные, не входящие в вывод, и он останется выводом,
но уже без новых констант. Лемма доказана.

На самом деле эта лемма верна для произвольного расширения
сигнатуры (можно добавлять не только константы, но и
функциональные символы любой валентности, а также предикатные
символы). Чтобы удалить новые символы из вывода, поступаем так.
Все термы вида~$f(\ldots)$, где $f$\т добавленный функциональный
символ, мы заменяем на новую переменную (можно взять одну и ту
же переменную для всех новых символов и всех их вхождений). Все
атомарные формулы с новыми предикатными символами заменяем на
какую-либо замкнутую формулу (одну и ту же; какая именно
формула, роли не~играет).

\begin{problem}
        %
Проведите это рассуждение подробно.
        %
\end{problem}

Таким образом, мы можем говорить о выводимости формулы, не
уточняя, в какой именно сигнатуре (содержащей все использованные
в формуле предикатные и функциональные символы) мы ищем её
вывод.

Если принять теорему о полноте, по которой выводимость
равносильна общезначимости, независимость выводимости от
сигнатуры становится очевидной: истинность формулы не зависит от
интерпретации символов, которые в неё не входят. (Если
интерпретировать отсутствующие в формуле символы как постоянные
функции и предикаты, мы приходим к синтаксическому рассуждению,
упомянутому выше.)
